% ID: SEL-D0062-OA-E03
% ESTADO: validada. S062.
Sean $M$ «es mujer» y $A$ «aprueba». Los varones representan $0{,}41$ y las mujeres $0{,}59$. Las probabilidades conjuntas de aprobar son $0{,}41\cdot0{,}70=0{,}287$ para los varones y $0{,}59\cdot0{,}60=0{,}354$ para las mujeres.

\textbf{a)} Por probabilidad total, $P(A)=0{,}641$. Por Bayes, $P(M\mid A)=0{,}354/0{,}641\approx\boxed{0{,}5523}$.

\textbf{b)} $P(A^c)=0{,}359$ y $P(M\cap A^c)=0{,}59\cdot0{,}40=0{,}236$. Así, $P(M\mid A^c)=0{,}236/0{,}359\approx\boxed{0{,}6574}$.

\textbf{c)} Ambos tienen razón: en cada grupo, aprobados o suspensos, la proporción de mujeres supera $1/2$.
